\chapter{Conclusion}\label{ch:Conclusion}
\section{Conclusion}
As active voltage source rectifiers (VSRs) become increasingly critical in modern power applications—such
as EV charging stations and renewable energy integration—their reliability under non-ideal grid conditions
is of paramount importance. This thesis investigated the adverse effects of unbalanced grid voltages, which
inherently generate negative-sequence currents and induce detrimental 100-Hz power oscillations on the DC-link.\par
The core contribution of this research lies in the integration of Symmetrical Component Extraction with an
Model Predictive Control (MPC) algorithm. By mapping the coupled three-phase system into decoupled
positive and negative sequence networks, an accurate discrete-time prediction model was formulated.
The controller was designed to regulate the square of the DC-link voltage (energy) to linearize the
power dynamics. A carefully structured multi-objective cost function was developed to simultaneously
track the DC energy reference, enforce unity power factor, and suppress second-order active and reactive
power ripples.\par
The effectiveness of the proposed architecture was rigorously tested against realistic grid fault scenario
(Math Bollen's Type B, F, and G sags) using a switching nonlinear model in Simulink. The simulation results
conclusively demonstrated that the proposed controller heavily attenuates the negative-sequence currents and
stabilizes the DC output. The dynamic response was exceptionally smooth, achieving settling times ranging
from 127.8 ms to 140.5 ms and limiting voltage overshoots to a maximum of 1.55\% even under the complex
phase-shifting conditions of Type F and Type G faults. Ultimately, this thesis proves that predictive control,
when combined with accurate sequence decomposition, offers a highly resilient and fast-acting solution for
VSRs facing grid asymmetries.\par
\section{Future Work}
The findings of this thesis lay a strong foundation for advanced power converter control, opening several
promising directions for future investigation:
\begin{itemize}
    \item \textbf{Predictive Model Update:} Future research could explore adaptive MPC frameworks that update
    the prediction model parameters in real-time to account for system nonlinearities, parameter variations,
    and unmodeled dynamics, further enhancing robustness and performance. The influence of nonlinearites can be 
    seen in the simulation results, where the nonlinear switching model shows slightly higher overshoots and 
    settling times compared to the linear model.
    \item \textbf{Grid Voltage Sensorless Control:} To reduce hardware costs and increase system reliability,
    future iterations of this controller could incorporate advanced observers (such as Extended Kalman Filters) to
    estimate the grid voltages and sequence components, eliminating the need for physical AC voltage sensors.
    \item \textbf{Experimental Validation:} Implementing the proposed control strategy on a hardware testbed 
    to validate its performance under real-world conditions, including measurement noise and component non-idealities.
\end{itemize}